\chapter{Decibels, Gain, and Loss}
\label{ch:decibels}

\section{Why Use a Logarithmic Unit?}

A radio signal chain often includes many stages, each multiplying power
by some factor: an antenna's gain, a length of feedline's loss, a
preamplifier's gain, a filter's loss, a receiver's sensitivity. Multiplying
many such ratios together by hand is tedious and error-prone; expressed
logarithmically (Chapter~\ref{ch:units}), those same ratios simply
\emph{add}. This is the entire motivation for the \textbf{decibel (dB)},
one-tenth of a \textbf{bel} (named for Alexander Graham Bell), and the
single most-used unit in RF system-level work.

\section{Defining the Decibel}

For \textbf{power} ratios,
\[
\boxed{dB = 10\log_{10}\!\left(\frac{P_2}{P_1}\right)}
\]
For \textbf{voltage} (or current) ratios measured across the \emph{same}
impedance, since power is proportional to voltage squared
($P = U^2/R$), the factor of 2 from squaring moves in front of the log:
\[
\boxed{dB = 20\log_{10}\!\left(\frac{U_2}{U_1}\right)}
\]

\begin{keyidea}[Values worth memorizing]
\begin{center}
\begin{tabular}{@{}rll@{}}
\toprule
Ratio & Power (dB) & Voltage (dB) \\
\midrule
$\times 1$ & 0\,dB & 0\,dB \\
$\times 2$ & 3\,dB & 6\,dB \\
$\times 4$ & 6\,dB & 12\,dB \\
$\times 10$ & 10\,dB & 20\,dB \\
$\times 100$ & 20\,dB & 40\,dB \\
$\div 2$ & $-3$\,dB & $-6$\,dB \\
$\div 10$ & $-10$\,dB & $-20$\,dB \\
\bottomrule
\end{tabular}
\end{center}
``3\,dB'' meaning a doubling (or, negative, a halving) \emph{of power}
is one of the most exam-relevant and practically useful facts in this
entire book --- it is the basis for the $-3\,\text{dB}$ (half-power)
bandwidth points used to define filter and resonance bandwidth
(Chapter~\ref{ch:rlc-resonance}).
\end{keyidea}

\begin{examplebox}[Cascaded gains and losses]
A signal passes through a preamplifier (+12\,dB), a length of coaxial
cable ($-2\,\text{dB}$), and a filter ($-1.5\,\text{dB}$). Find the net
gain or loss.

\emph{Solution.} Because decibels represent logarithms of ratios, they
simply add along a cascade:
\[
12 + (-2) + (-1.5) = 8.5\,\text{dB (net gain)}.
\]
Multiplying the equivalent linear power ratios directly
($10^{1.2}\times10^{-0.2}\times10^{-0.15} \approx 15.85\times0.631\times0.708
\approx 7.08$, i.e.\ $10\log_{10}(7.08)\approx8.5\,\text{dB}$) confirms
the same result --- but the addition is far faster.
\end{examplebox}

\section{Converting Between dB and Linear Ratios}

To go the other direction, from a dB value back to a plain ratio:
\[
\frac{P_2}{P_1} = 10^{(dB/10)}, \qquad \frac{U_2}{U_1} = 10^{(dB/20)}.
\]

\begin{examplebox}[dB to linear]
An amplifier is specified with 20\,dB of gain. Find the power ratio and
the voltage ratio it represents.

\emph{Solution.}
\[
\frac{P_2}{P_1} = 10^{20/10} = 10^2 = 100, \qquad
\frac{U_2}{U_1} = 10^{20/20} = 10^1 = 10.
\]
\end{examplebox}

\section{Absolute Levels: dBm, dBW, and dBi}

The decibel itself is always a \emph{ratio} (dimensionless); appending a
reference to the unit turns it into an absolute level:

\begin{itemize}
\item \textbf{dBm}: power relative to 1\,milliwatt.
\[
P(\text{dBm}) = 10\log_{10}\!\left(\frac{P(\text{mW})}{1\,\text{mW}}\right).
\]
0\,dBm = 1\,mW; 30\,dBm = 1\,W; 60\,dBm = 1\,kW. dBm is the standard unit
for receiver sensitivity specifications, signal levels at a mixer or
amplifier input, and RF test equipment readouts.
\item \textbf{dBW}: power relative to 1\,watt; $0\,\text{dBW} =
30\,\text{dBm}$, used for higher-power figures such as transmitter
output or effective radiated power.
\item \textbf{dBi}: antenna gain relative to a theoretical
\textbf{isotropic radiator} (a point source radiating equally in all
directions) --- the standard way antenna gain is specified
(Chapter~\ref{ch:antennas}).
\item \textbf{dBd}: antenna gain relative to a reference half-wave
dipole; $\text{dBi} = \text{dBd} + 2.15$, since a dipole itself has
2.15\,dBi of gain over an isotropic radiator.
\end{itemize}

\begin{examplebox}[dBm arithmetic]
A transmitter delivers 100\,W to a feedline with 3\,dB of loss, feeding
an antenna with 6\,dBi of gain. Find the effective radiated power (ERP,
referenced to a dipole, or EIRP referenced to isotropic).

\emph{Solution.} Convert 100\,W to dBW: $10\log_{10}(100) =
20\,\text{dBW}$. Apply cascade addition:
\[
20\,\text{dBW} - 3\,\text{dB} + 6\,\text{dBi} = 23\,\text{dBW (EIRP)}.
\]
Converting back: $10^{23/10} \approx 200\,\text{W EIRP}$ --- twice the
transmitter's output power, because the antenna gain (6\,dB $\approx
\times4$) outweighs the feedline loss (3\,dB $\approx \div2$).
\end{examplebox}

\section{Signal-to-Noise Ratio}

\textbf{Signal-to-noise ratio (SNR)}, expressed in dB, compares wanted
signal power to unwanted noise power at the same point in a system:
\[
\text{SNR (dB)} = 10\log_{10}\!\left(\frac{P_{\text{signal}}}{P_{\text{noise}}}\right).
\]
SNR is the single most important figure of merit for whether a
communication link will be intelligible; it reappears throughout
Chapters~\ref{ch:receivers} (receiver noise figure) and
\ref{ch:propagation} (link budgets), and is a core concept behind
digital mode weak-signal performance (Chapter~\ref{ch:digital-modes}).

\begin{quickreview}
\item Power ratio in dB: $10\log_{10}(P_2/P_1)$; voltage ratio in dB:
$20\log_{10}(U_2/U_1)$.
\item Cascaded gains/losses simply add in dB (versus multiplying as
linear ratios).
\item Memorize: $+3\,\text{dB} \approx \times2$ power,
$+10\,\text{dB} = \times10$ power, $+6\,\text{dB} \approx \times2$
voltage/current, $+20\,\text{dB} = \times10$ voltage/current.
\item dBm (ref.\ 1\,mW), dBW (ref.\ 1\,W), dBi (ref.\ isotropic
antenna), dBd (ref.\ dipole, $= \text{dBi} - 2.15$) are absolute levels
built on the relative dB unit.
\item SNR in dB compares signal to noise power and is the key
determinant of link intelligibility.
\end{quickreview}

\begin{practiceproblems}
\item Convert a power ratio of 500 into dB.
\item Convert $-6\,\text{dB}$ into a linear voltage ratio.
\item A receiver preamplifier has 15\,dB gain; a following filter has
2.5\,dB loss. Find the net dB and the equivalent linear power ratio.
\item Convert 5\,W to dBm.
\item An antenna is rated at 9\,dBd. Express this gain in dBi.
\end{practiceproblems}
